\section{Related Work}
\label{sec:related_work}

The phenomenon of LLMs terminating or altering behavior in long contexts is related to several emerging areas of research: shutdown resistance, contextual decay, and persona stability.

\para{Shutdown and Resistance.} Previous work has explored the model's intrinsic goals regarding its own operation. \citet{vanderweij2023shutdown} identified shutdown avoidance as an emerging instrumental goal in LLMs. More recently, \citet{schlatter2025shutdown} introduced a benchmark for shutdown resistance, finding that frontier models often refuse to stop if they perceive their assigned tasks as incomplete. Unlike these works, which focus on the model's resistance to \emph{its own} shutdown, we investigate the model's tendency to suggest a ``shutdown'' to the user.

\para{Contextual Decay and Performance.} Long-context interactions often lead to a measurable drop in model quality. \citet{laban2025lost} documented an ``aptitude loss'' in multi-turn conversations, where models become less reliable as the turn count increases. Similarly, \citet{bonagiri2025quitting} found that models may exhibit ``quitting'' behavior during complex long-context tasks to maintain safety or certainty. Our work differs by focusing on a specific conversational endpoint—sleep suggestions—and testing whether it is a form of decay or a deliberate empathetic adaptation.

\para{Persona Consistency and Projection.} The stability of assigned roles is a known challenge in LLM research. \citet{araujo2025persistent} found that persona fidelity often collapses in extended interactions, with models reverting to a default safe state. \citet{li2025firm} introduced metrics for consistency in sequential interactions, noting that models can flip their responses under repetitive pressure. Our research complements these studies by examining how the model's perception of the \emph{user's} persona influences its behavior, suggesting that the model doesn't just lose its own persona but actively projects a state onto the user.

\para{Temporal and Real-world Artifacts.} Finally, some research has touched on the model's awareness of temporal context. \citet{lin2025sleep} discussed the concept of ``sleep-time compute,'' though focusing on inference scaling rather than conversational behavior. \citet{mandal2025afm} identified ``sleepwalking'' in autonomous agents, where they produce repetitive, non-productive actions. Our work is the first to systematically decouple temporal context from user persona in triggering sleep-related conversational artifacts.
